
\chapter{Introdução}
\label{chap:Introdução}

Este projeto propõe uma nova abordagem baseada em sistemas neuro-fuzzy evolutivos para detecção e prevenção de fraudes em operações do setor agropecuário, utilizando dados reais de Guias de Trânsito Animal (GTA). A proposta explora a integração entre redes neurais artificiais e sistemas fuzzy evolutivos, permitindo que o modelo aprenda automaticamente funções de pertinência e regras de decisão a partir dos dados, mesmo diante de incertezas, ruídos e padrões em constante transformação. Essa arquitetura neuro-fuzzy evolutiva é especialmente adequada para ambientes dinâmicos, pois combina a capacidade de adaptação e generalização das redes neurais com a interpretabilidade e flexibilidade dos sistemas fuzzy, tornando-se robusta frente à evolução dos padrões de fraude.

O contexto do setor agropecuário brasileiro, marcado por grandes volumes de dados e vulnerabilidade a práticas fraudulentas, exige soluções inteligentes e adaptativas. A utilização de dados reais de GTAs, que registram movimentações de animais e são fundamentais para o controle sanitário e fiscal, possibilita a criação de um sistema capaz de identificar irregularidades e prevenir fraudes, como as observadas em operações citadas nesse como CIRA, Rei do Gado e Cacus.

O trabalho está estruturado para abordar desde a fundamentação teórica dos sistemas neuro-fuzzy evolutivos, passando pela análise e preparação dos dados, até a implementação, validação e discussão dos resultados obtidos. Espera-se contribuir para o avanço científico e prático na detecção de fraudes em ambientes complexos e dinâmicos, fornecendo suporte efetivo à fiscalização, à tomada de decisão e à integridade do setor agropecuário.

Os órgãos de defesa agropecuária são responsáveis pela coordenação das atividades de defesa agropecuária, com foco na segurança alimentar, saúde pública e qualidade dos produtos agropecuários. Essas instituições desempenham funções como fiscalização sanitária, controle de doenças e pragas, rastreabilidade de produtos e apoio técnico ao setor agropecuário. Muitas dessas organizações utilizam sistemas de gerenciamento do setor agropecuário que facilitam o controle sanitário, rastreabilidade e monitoramento de doenças, promovendo maior transparência e eficiência na fiscalização, além de fortalecer a competitividade do setor agropecuário. 

No sistema de gerenciamento do setor agropecuário, são coletadas diversas informações, criando um vasto banco de dados que já foi utilizado em trabalhos acadêmicos. Esses sistemas geram volumes significativos de dados, tratando-se de uma base de dados crescente que no momento conta com uma estrutura de aproximadamente 8.298.490 registros de Guias de Trânsito Animal (GTA) apenas para movimentações internas em determinadas regiões. Esses dados têm sido empregados em diversas pesquisas, como as que exploram a movimentação de bovinos, com base nas GTAs.

Nesse aspecto, o estudo de \citeonline{costa2023} analisa as redes de movimentação de bovinos entre municípios de Minas Gerais, utilizando dados da GTA coletados entre 2013 e 2023. A pesquisa identificou padrões de movimentação concentrados na região do Triângulo Mineiro e Alto Paranaíba, classificando a rede como altamente conectada, com características típicas de grafos do tipo small-world. 

Outro trabalho que pode ser citado é o artigo de \citeonline{azevedo2022}, que caracteriza a rede de movimentação de bovinos no estado de Goiás entre 2010 e 2016, utilizando dados de 4.697.239 Guias de Trânsito Animal (GTA). Os dados foram submetidos a análises estatísticas descritivas e de redes, revelando um intenso fluxo de animais, especialmente nas regiões noroeste e sul do estado, com destinos voltados para engorda, recria e abate. 

Esses dois estudos demonstram como a aplicação de análises de redes e outras ferramentas tecnológicas pode contribuir para melhorar a eficiência e a eficácia das estratégias de controle sanitário e epidemiológico, promovendo a sustentabilidade e a competitividade do setor agropecuário brasileiro.

O setor agropecuário brasileiro, especialmente o segmento de gado bovino, tem se mostrado vulnerável a práticas fraudulentas, como evidenciado em diversas operações de combate à sonegação fiscal e fraudes no comércio de gado. Podemos citar alguns exemplos de deflagrações fraudulentas encontradas na íntegra, como a:

\begin{itemize}
\item Operação do CIRA, deflagrada pelo Ministério Público de Minas Gerais em 2022, que, ao longo de 11 meses, emitiu 21 notas fiscais Guias de Trânsito Animal em nome do produtor rural sem que ele soubesse, totalizando uma quantia de R\$ 1.670.451,50. 
\item A Operação Rei do Gado, conduzida pela Receita Federal, investigou fraudes fiscais de R\$ 1,4 bilhão em transações ilegais envolvendo a venda de bovinos \cite{cnn2023}. 
\item A Operação Boi na Linha \cite{g12016} também destacou as dificuldades de fiscalização no setor, com a prisão de envolvidos em fraudes envolvendo a venda e o transporte ilegal de gado. 
\item Por fim, a recente Operação Cacus, realizada pela Polícia Civil de Minas Gerais, investigou servidores estaduais suspeitos de estarem envolvidos em um esquema de movimentação ilegal de 83 mil cabeças de gado \cite{g12024}.
\end{itemize}

Essas operações exemplificam a necessidade urgente de implementação de tecnologias inteligentes para monitorar, identificar e prevenir fraudes nesse setor. A utilização de modelos baseados em sistemas inteligentes, como o que se propõe neste estudo, poderá contribuir significativamente para reduzir as irregularidades, melhorar a fiscalização e garantir maior integridade ao sistema de emissão de documentos e transações de bovinos.

Para tanto, podemos também exemplificar alguns trabalhos existentes que investigam fraudes em diversos setores, utilizando diferentes tipos de ferramentas e propondo novos paradigmas para predizer possíveis fraudes no sistema. Alguns desses trabalhos são citados a seguir. 

O estudo de \citeonline{bao2019} propõe um modelo de previsão de fraudes contábeis utilizando técnicas de sistemas inteligentes, destacando a abordagem de ensemble learning. Em vez de utilizar razões financeiras tradicionais, os autores empregam dados financeiros brutos como entrada para o modelo que também introduziu métricas de avaliação, como o Normalized Discounted Cumulative Gain (NDCG), para priorizar casos de fraude com maior probabilidade de ocorrência, tornando-o prático para reguladores e auditores com recursos limitados. 

O trabalho de \citeonline{singh2019} apresenta um modelo de detecção de fraudes em cartões de crédito baseado em Modelos Ocultos de Markov (HMM). A proposta envolve treinar o HMM com o comportamento normal dos usuários em transações financeiras, possibilitando identificar fraudes por meio de desvios nos padrões de gasto. O modelo é dividido em duas fases: treinamento, onde os dados históricos do usuário são utilizados para criar um perfil transacional, e detecção, que analisa novas transações em tempo real. 

O artigo de \citeonline{coelho2006} apresenta o FControl®, um sistema computacional híbrido baseado em inteligência computacional, projetado para detecção de fraudes em transações de comércio eletrônico com cartões de crédito. O FControl® utiliza uma abordagem híbrida de inteligência computacional baseada em sistema neuro-nebuloso (neuro-fuzzy) e otimização por estratégia evolutiva com mecanismos de auto-adaptação de parâmetros. A prevenção de fraude em cartão de crédito representa uma importante aplicação comercial para métodos de previsão e inteligência computacional, combinando sistemas neuro-nebulosos, redes neurais artificiais, computação evolutiva e estratégias de otimização. O sistema se destaca pela rapidez e eficiência na detecção de fraudes, demonstrando a eficácia de abordagens híbridas que integram diferentes paradigmas de sistemas inteligentes para problemas de classificação e detecção de anomalias em tempo real. 

A pesquisa de \citeonline{mallela2024} explora o uso de técnicas de Machine Learning (ML) para aprimorar os mecanismos de detecção de fraudes em instituições financeiras. A pesquisa avalia modelos supervisionados, não supervisionados e de aprendizado por reforço, destacando a importância da engenharia de características e do pré-processamento de dados para a robustez dos modelos. 

No trabalho apresentado por \citeonline{hooda2018}, que apresenta uma abordagem para classificação de empresas fraudulentas utilizando algoritmos de sistemas inteligentes em auditorias externas. O estudo utilizou dados de 777 empresas de 14 setores diferentes, aplicando o algoritmo Particle Swarm Optimization (PSO) para seleção de características relevantes e comparando dez modelos de classificação. 

\citeonline{zhou2023} apresentaram uma abordagem inovadora para detecção de fraudes em demonstrações financeiras utilizando um grafo de conhecimento de duas camadas (FSFD-TLKG). A camada semântica do grafo modela os relacionamentos de subordinação entre variáveis financeiras, enquanto a camada sintática explora as relações de articulação inerentes às demonstrações financeiras. 

Esses estudos demonstram novos e combinação de diferentes métodos, como aprendizado supervisionado, não supervisionado, aprendizado por reforço, IA e grafos de conhecimento, permite a criação de sistemas robustos, eficientes e dinâmicos. Além disso, a aplicação dessas tecnologias em setores como auditoria, comércio eletrônico e finanças pode melhorar substancialmente a capacidade de identificar fraudes e reduzir os prejuízos associados. O desenvolvimento de soluções preditivas e automatizadas para detecção de fraudes está se tornando cada vez mais relevante, alinhando-se com as necessidades de eficiência e segurança nos sistemas modernos.

No contexto específico de sistemas evolutivos para detecção de fraudes, a incorporação de métodos de pertinência fuzzy emerge como uma abordagem promissora. O artigo de \citeonline{coelho2006}, por exemplo, apresenta o FControl®, um sistema híbrido neuro-fuzzy para detecção de fraudes. Os sistemas fuzzy evolutivos combinam a capacidade de lidar com incertezas e imprecisões inerentes aos dados agropecuários, por meio de funções de pertinência, com a adaptabilidade estrutural dos algoritmos evolutivos. Essa combinação permite que o sistema modele adequadamente as nuances dos padrões fraudulentos, que frequentemente apresentam características ambíguas e limites não bem definidos.

Os sistemas neuro-fuzzy evolutivos, pioneiramente desenvolvidos por Kasabov e Song \cite{kasabov-2002} com o DENFIS (Dynamic Evolving Neural-Fuzzy Inference System), representam uma evolução natural desta abordagem. Angelov e Filev \cite{angelov-2004b} expandiram esta linha de pesquisa com métodos de identificação online de modelos Takagi-Sugeno, enquanto Lughofer \cite{Lughofer2011} consolidou as metodologias e conceitos avançados desses sistemas evolutivos. Mais recentemente, trabalhos como o de Skrjanc et al. \cite{Skrjanc2019a} demonstraram a eficácia de classificadores evolutivos em ambientes de streaming de dados, características essenciais para aplicações em tempo real.

Portanto, este trabalho de pesquisa propõe o desenvolvimento de um modelo inteligente para detecção e prevenção de fraudes em sistemas de gestão, com validação por meio de dados gerados por sistemas de gerenciamento do setor agropecuário, a partir de um estudo de caso real. A pesquisa atende a uma necessidade do setor de criar métodos para mitigar fraudes no sistema. A proposta se destaca por sua abordagem específica, adaptada às particularidades do sistema agropecuário em questão, a partir de análise dos dados coletados ao longo dos anos fornecerá insights valiosos para a criação de um sistema inteligente capaz de detectar irregularidades e prevenir fraudes, como as observadas nas operações CIRA, Rei do Gado e Cacus.

O projeto está organizado da seguinte maneira: na Seção \ref{chap:Problema} é apresentada uma caracterização do problema; na Seção \ref{chap:objetivos} são definidos os objetivos gerais e específicos; na Seção \ref{chap:metodologia} descreve-se a metodologia que será utilizada no desenvolvimento da pesquisa; e na Seção \ref{chap:cronograma} são apresentados o cronograma de execução e a formação acadêmica realizada.
